% Consolidated guide-selftest04 from 97b62cc61307eba158a4498b55de3a5fdef3ff43; source fragments preserved.

% === BEGIN guide-selftest04-choices.tex blob 64b98f7ff8ad88542a616b5ccf2b4f4e5dd8f0bf ===
% 第16批：扫描原书自测题4选择，共10题；题意独立转述。
\subsection{自测题4：简谐振动、行波与干涉选择题}

\sourcetest{4.1.1}{93}{101}
\problemtext 简谐振动 $x=A\cos(\omega t+\pi/4)$；求 $t=T/4$ 时的加速度（原题选项）。
\workings 由 $a=-\omega^2x$，此时相位 $3\pi/4$，故 $\boxed{a=+\sqrt2 A\omega^2/2}$，选 \textbf{B}。
\crosscheck 加速度永远指向平衡位置，不能把速度公式误作加速度。

\sourcetest{4.1.2}{93}{101}
\problemtext 在简谐振动端点和平衡位置判断速度、加速度的性质。
\workings 平衡点有 $|v|=A\omega$ 与 $a=0$；端点有 $v=0$ 与 $|a|=A\omega^2$。因此题中只有“平衡位置正向运动，速度最大、加速度为零”成立，选 \textbf{C}。
\crosscheck “速度最大”指有向速度取得正最大值，若逆向经过则为负最大绝对值。

\sourcetest{4.1.3}{93}{101}
\problemtext 题4.1.1图中 $v(0)=v_{\max}/2$ 且曲线初始上升，求余弦位移式的初相。
\workings 若 $x=A\cos(\omega t+\varphi)$，则 $v=-A\omega\sin(\omega t+\varphi)$。$\sin\varphi=-1/2$ 且 $dv/dt=a(0)>0$，故 $\cos\varphi<0$，得到 $\boxed{\varphi=-5\pi/6\pmod{2\pi}}$，选 \textbf{C}。
\crosscheck 图中纵坐标为速度而非位移，初相判定必须兼顾速度图的斜率。

\sourcetest{4.1.4}{93}{101}
\problemtext 已知 $x=0.04\cos(2\pi t+\pi/3)$（SI），从 $t=0$ 起首次到达 $x=-2$ cm 且朝 $+x$ 方向运动的时间。
\workings 目标相位需满足 $\cos\theta=-1/2$、$\sin\theta<0$，故 $\theta=4\pi/3$；从 $\pi/3$ 走过 $\pi$，$\Delta t=\pi/(2\pi)=\boxed{1/2\ {\rm s}}$，选 \textbf{C}。
\crosscheck 必须同时利用位移和运动方向，另一个根 $2\pi/3$ 对应反向。

\sourcetest{4.1.5}{93--94}{NaN}
\problemtext 振幅 $4$ cm、周期 $2$ s；$t=0$ 第一次经过 $x=-2$ cm 时沿负向，求第二次经过该位置的时刻。
\workings 首经过可取相位 $2\pi/3$（$\sin\varphi>0$），第二次对应 $4\pi/3$，故 $\Delta t=(2\pi/3)/\omega=(2\pi/3)/\pi=\boxed{2/3\ {\rm s}}$，选 \textbf{B}。
\crosscheck “第二次”指下一次经过同一位置，方向变为正向，不是一整个周期。

\sourcetest{4.1.6}{94}{102}
\problemtext 波动式 $y=0.1\cos(3\pi t-\pi x+\pi)$（SI）；根据题4.1.2图辨别振幅、波长、两点相位差和波速的正确表述。
\workings 振幅 $A=0.1$ m、波长 $\lambda=2\pi/\pi=2$ m，传播速率 $u=3$ m/s。图中 $a,b$ 相隔波长四分之一，故相位差的大小 $\pi/2$，选 \textbf{C}。
\crosscheck 振幅不取负值；相位差绝对值与有向相位差分开说明。

\sourcetest{4.1.7}{94}{102}
\problemtext 题4.1.3图给出向 $-x$ 传播的横波某瞬时波形，判断 A、B、C、D 点的横向振动速度符号。
\workings 对左行波 $y=f(x+ut)$，有 $\partial y/\partial t=u\,\partial y/\partial x$。图中 D 点波形斜率小于零，故 $v_{y,D}<0$；其余陈述不符，选 \textbf{D}。
\crosscheck 横向振动速度不同于沿 $x$ 方向的波速，方向由传播方向与局部斜率共同决定。

\sourcetest{4.1.8}{94}{102}
\problemtext 行进的简谐弹性波中某质点恰处于平衡位置，比较该处动能和势能。
\workings 行波位移 $y=A\cos(\omega t-kx)$；质点速度 $y_t=-A\omega\sin\theta$，局部应变 $y_x=Ak\sin\theta$。两种能量密度均与 $\sin^2\theta$ 成正比；$y=0$ 时二者同时最大，选 \textbf{C}。
\crosscheck 不能机械套用单个弹簧振子的“平衡点势能零”，行波弹性势能由空间梯度决定。

\sourcetest{4.1.9}{94}{102}
\problemtext 两相干源到点 $P$ 的距离分别为 $r_1,r_2$，初相 $\varphi_1,\varphi_2$；求 $P$ 的干涉极大相位条件。
\workings 到达 $P$ 的两相位为 $\varphi_j-2\pi r_j/\lambda$，所以 $\boxed{\varphi_2-\varphi_1+2\pi(r_1-r_2)/\lambda=2k\pi}$；选 \textbf{D}。
\crosscheck 此式允许两源初相不相等；单靠路程差为整数波长并不足够。

\sourcetest{4.1.10}{94}{102}
\problemtext 等振幅、同频率的两列反向传播简谐波 $y_{1,2}=A\cos[2\pi(\nu t\mp x/\lambda)]$ 叠加，求驻波的波节。
\workings 由和差化积 $y=2A\cos(2\pi\nu t)\cos(2\pi x/\lambda)$。波节满足 $\cos(2\pi x/\lambda)=0$，故 $\boxed{x=(2k+1)\lambda/4,\ k\in\mathbb Z}$，选 \textbf{D}。
\crosscheck 波节与波腹相隔 $\lambda/4$，相邻波节相隔 $\lambda/2$。
% === END guide-selftest04-choices.tex ===

% === BEGIN guide-selftest04-blanks.tex blob 5f63e7769f58151bcad74969ea2b80d529901c60 ===
% 第16批：原书印刷第94—95页；源图题标注复核条件。
\subsection{自测题4：简谐振动、波形及驻波填空}

\sourcetest{4.2.1}{94}{102}
\problemtext 频率 $0.25$Hz，$t=0$ 时位移为 $-0.37$cm 且速度为零；求振幅及余弦型振动式。
\workings 零速即转折点，$A=0.37$cm，$\omega=2\pi f=\pi/2$rad/s。故 $\boxed{x=0.37\cos(\pi t/2+\pi)\ {\rm cm}}$。
\crosscheck 初始在负端点，必须取初相 $\pi$，不是0。

\sourcetest{4.2.2}{95}{103}
\problemtext 竖直弹簧振子位移为振幅的一半时，求动能与总能量的比，以及由静伸长 $\Delta l$ 表示的周期。
\workings $E=\frac12kA^2$；位移为 $A/2$ 时 $K/E=1-(1/2)^2=\boxed{3/4}$。静平衡 $k\Delta l=mg$，$T=2\pi\sqrt{m/k}=\boxed{2\pi\sqrt{\Delta l/g}}$。
\crosscheck 势能零点取弹簧振子的重力修正平衡位置。

\sourcetest{4.2.3}{95}{103}
\problemtext 题4.2.1图示振幅为4个纵轴单位，$t=0$ 时 $x=-2$ 且正在向正向运动，$t=2$s 下降经过平衡点；写出周期和余弦式初相。
\workings 下降过平衡点时的相位为 $\pi/2$（模 $2\pi$）；$t=0$ 为 $\cos\varphi=-1/2$ 且 $v>0$，可取 $\varphi=-2\pi/3$。从 $-2\pi/3$ 到 $\pi/2$ 的相位跨度为 $7\pi/6$，据此 $\omega=7\pi/12$rad/s、$\boxed{T=24/7\ {\rm s}}$。
\crosscheck 图中零点的时间标注需要独立对照：若印刷图仅为示意而非精确刻度，周期不应由图形长度估读；本条保留待复核。

\sourcetest{4.2.4}{95}{103}
\problemtext 同向同频简谐振动合振幅为20cm，第二振动振幅 $10\sqrt3$cm，且 $\varphi_2-\varphi_1=\pi/6$；求第一振幅及逆序相位差。
\workings $20^2=A_1^2+(10\sqrt3)^2+2A_1(10\sqrt3)\cos(\pi/6)$，得 $A_1^2+30A_1-100=0$。取正根 $\boxed{A_1=5\sqrt{13}-15\simeq3.028\ {\rm cm}}$；$\boxed{\varphi_1-\varphi_2=-\pi/6}$。
\crosscheck 振幅必须取非负根；相位差反序改变符号。

\sourcetest{4.2.5}{95}{103}
\problemtext 相干波源 $S_1,S_2$ 相距 $3\lambda/4$，两源等振幅；在 $S_1$ 左侧各点合波强度为单源4倍，求两源初相差。
\workings 该侧任意点的传播路程差为 $r_2-r_1=3\lambda/4$；强度达到4倍表明同相抵达：$\varphi_2-\varphi_1-2\pi(3/4)=2k\pi$，故 $\boxed{\varphi_2-\varphi_1=3\pi/2\equiv-\pi/2\pmod{2\pi}}$。
\crosscheck “某侧各点”成立的原因是两源在同一直线上，路程差恒定。

\sourcetest{4.2.6}{95}{103}
\problemtext 题4.2.3图给出同一右行波在 $t=0$ 与 $T/4$ 的波形；画出标记点 $P$ 的振动图线。
\workings 从图上 $P$ 在 $t=0$ 经过平衡位置，且波形局部斜率为负；右行波满足 $y_t=-u y_x$，故此时 $v_y>0$。所以可写 $\boxed{y_P(t)=A\sin(2\pi t/T)}$。振动图线从原点以正斜率上升，$T/4$ 达到 $+A$。
\crosscheck 由两幅图的峰值右移 $\lambda/4$ 确认传播方向；与原图 P 的下降空间斜率对应上升时间斜率。

\sourcetest{4.2.7}{95}{103}
\problemtext 题4.2.4图示右行波在 $t=T/4$ 的空间曲线，波速330m/s，振幅0.10m；横轴两相邻波谷距离4m。求波动方程。
\workings $\lambda=4$m，$k=\pi/2$rad/m，$\omega=2\pi(330/4)=165\pi$rad/s。图中该时刻 $x=0$ 过零且波形斜率负，故可选 $\boxed{y=0.10\cos(165\pi t-\pi x/2+\pi)}$（SI）。
\crosscheck 代入 $t=T/4=\pi/(2\omega)$，得 $y(x,T/4)=-0.10\sin(\pi x/2)$，波谷在 $x=1$m，符合原图。

\sourcetest{4.2.8}{95}{103}
\problemtext 右行波 $y=A\cos[2\pi(t/T-x/\lambda)+\varphi]$，分别求 $x=-\lambda$ 质点振动式及以旧坐标 $x=\lambda$ 为原点、反向取轴的新波动式。
\workings 代入 $x=-\lambda$，$\boxed{y_{-\lambda}(t)=A\cos(2\pi t/T+\varphi)}$。若新坐标 $x'=\lambda-x$，则 $x=\lambda-x'$，所以 $\boxed{y=A\cos[2\pi(t/T+x'/\lambda)+\varphi]}$。
\crosscheck 空间坐标轴反向后，相位中空间项的号随之反转；相移一个波长不改变相位模 $2\pi$。

\sourcetest{4.2.9}{95}{103}
\problemtext 左行入射波 $y_i=A\cos[2\pi(t/T+x/\lambda)]$ 在固定端 $x=0$ 发生等幅反射；求反射波式及 $x=2\lambda/3$ 处合振幅。
\workings 固定端位移相消，$\boxed{y_r=-A\cos[2\pi(t/T-x/\lambda)]}$。合波 $y=-2A\sin(2\pi t/T)\sin(2\pi x/\lambda)$，于 $x=2\lambda/3$ 得振幅 $\boxed{A_P=\sqrt3 A}$。
\crosscheck 固定端反射发生 $\pi$ 相变；反射波不是与入射波同向传播。

\sourcetest{4.2.10}{95}{103}
\problemtext 汽车以20m/s速度驶近前方观察者，鸣笛频率1000Hz，空气声速330m/s；求车前方声波波长。
\workings 移动声源朝前发出的相邻波峰间距 $\lambda=(c-v_s)/f=(330-20)/1000$，故 $\boxed{\lambda=0.310\ {\rm m}}$。
\crosscheck 这里车前方受压缩的波长小于静止声源波长0.330m；原题给的1000Hz是鸣笛本征频率。
% === END guide-selftest04-blanks.tex ===

% === BEGIN guide-selftest04-solutions.tex blob 61bd808f8131b78a398d5fd20c1426d648e604df ===
% 第16批：独立推导、自查；依赖源图的题标明判读前提。
\subsection{自测题4：振动与波动计算、问答}

\sourcetest{4.3.1}{96}{104}
\problemtext 质点位移 $x=0.24\cos(\pi t/2+\pi/3)$（SI），求从 $t=0$ 到首次达到 $x=-0.12$m、且 $v<0$ 的时间。
\workings 角频率 $\omega=\pi/2$；末态满足 $\cos\theta=-1/2$ 且 $\sin\theta>0$，故最早取 $\theta=2\pi/3$。初相 $\pi/3$，于是 $\boxed{\Delta t=(2\pi/3-\pi/3)/(\pi/2)=2/3\ {\rm s}}$。
\crosscheck 旋转矢量沿正向转过 $\pi/3$；另一个位移根 $4\pi/3$ 对应 $v>0$。

\sourcetest{4.3.2}{96}{104}
\problemtext 质量2kg的物块放于竖直振动平板上，板周期0.5s、振幅4cm；板初始在上方最大位移处，求物块对板压力与开始离板的临界振幅。
\workings 取竖直向上为正、$z=A\cos\omega t$，$\omega=2\pi/T=4\pi$rad/s。接触时 $N-mg=m\ddot z=-m\omega^2 z$，即 $N=m(g-\omega^2 z)$。若板在最高点，$N_{\min}=m(g-\omega^2A)\simeq2(9.8-16\pi^2\cdot0.04)=\boxed{6.97\ {\rm N}}$；最低点 $N_{\max}=m(g+\omega^2A)\simeq\boxed{32.23\ {\rm N}}$。临界最高点 $N_{\min}=0$，$A_c=g/\omega^2\simeq\boxed{6.21\ {\rm cm}}$。
\crosscheck 原题表述“物体对平板的压力”数值等于平板支持力；从最高点出发处压力为最小值，尚未离板。

\sourcetest{4.3.3}{96}{104}
\problemtext 题4.3.1图给出向左传播的同一简谐波在 $t=0$ 和 $t=2$s 的波形，两图向左平移20m；已知原点处 $t=0$ 过零且初速度为正，$t=2$s 位移为 $A/\sqrt2$，求原点振动式和波动式。
\workings 左行波写为 $y=A\sin(\omega t+kx)$。在原点两时刻相位由0增加 $\pi/4$，故取图示最小相移 $\omega(2\,{\rm s})=\pi/4$，$\omega=\pi/8$rad/s。两图相移20m对应 $k\,20=\pi/4$，$k=\pi/80$rad/m，从而 $u=\omega/k=10$m/s。故 $\boxed{y(0,t)=A\sin(\pi t/8)}$，$\boxed{y(x,t)=A\sin(\pi t/8+\pi x/80)}$。
\crosscheck 波形在2s内向左移20m；若不排除整周期别名 $\omega 2=\pi/4+2n\pi$，单凭两幅快照不能唯一决定频率；本解采用原图绘示的最短相移分支，需保留图示条件标记。

\sourcetest{4.3.4}{96}{104}
\problemtext 一右行波速500m/s，点 $P$ 位于 $x=1$m，$y_P=0.03\cos(500\pi t-\pi/2)$（SI）；求全波波动式及 $t=0$ 波形。
\workings $\omega=500\pi$rad/s，$k=\omega/u=\pi$rad/m。设 $y=0.03\cos(500\pi t-\pi x+\varphi)$；于 $x=1$ 应有 $-\pi+\varphi=-\pi/2$（模 $2\pi$），故 $\varphi=\pi/2$。$\boxed{y(x,t)=0.03\cos(500\pi t-\pi x+\pi/2)}$；$\boxed{y(x,0)=0.03\sin(\pi x)}$。
\crosscheck 验证在 $x=1$m 时恢复题给 P 点振动式；波长2m。

\sourcetest{4.3.5}{96}{104}
\problemtext 右行声波频率 $f_n$、波长 $\lambda_n$，遇到向左匀速移动的刚性平面（速率 $v_0<c$）后反射；求 $\Delta t$ 内入射波峰数、反射波最前端距离以及反射波长。
\workings 入射波与反向运动平面的相遇速率为 $c+v_0$，碰到平面的波峰数 $\boxed{N=f_n(1+v_0/c)\Delta t}$。最初反射波峰左行速率 $c$，平面左移 $v_0\Delta t$，故两者末距 $\boxed{d=(c-v_0)\Delta t}$。反射波峰间距 $\lambda_r=(c-v_0)/(f_n(1+v_0/c))=\boxed{\lambda_n(c-v_0)/(c+v_0)}$。
\crosscheck 此题壁面靠近声源，反射波波长缩短；当 $v_0\to0$ 恢复 $\lambda_r=\lambda_n$。

\sourcetest{4.4.1}{96}{104}
\problemtext 题4.4.1图为向 $+x$ 传播的余弦式简谐波 $t=0$ 的波形；求图中 O、2、3 等标点的初相并说明如何画出 $t=T/4$ 波形。
\workings 图中 $t=0$ 的波形可写作 $y(x,0)=A\sin(kx)$（相邻过零点间距为半波长）；右行波为 $\boxed{y(x,t)=A\sin(kx-\omega t)}$。改写成余弦振动形式 $A\cos(\omega t-kx+\varphi)$，初相在 $O(x=0)$ 为 $\pi/2$，$x=2$（下降过零）为 $-\pi/2$，$x=3$（波谷）为 $-\pi$（均按 $2\pi$ 取模）。$t=T/4$ 波形是原图整体向右平移 $\boxed{\lambda/4}$。
\crosscheck 以题图横轴标注 $x=1,2,3,4$ 理解 $x$ 的单位；原图局部峰谷可通过代入波形验证。
% === END guide-selftest04-solutions.tex ===
